\documentclass[10pt]{amsart}
\usepackage[a4paper,margin=23mm]{geometry}
\usepackage{amsmath,amssymb,amsthm,mathtools,hyperref,graphicx,xspace,letltxmacro,xpatch}
\usepackage[english]{babel}
\usepackage{xfrac}

\usepackage{cases}
\usepackage{mathrsfs}
\let\mathcal\mathscr
\let\cal\mathcal
\makeatletter
\renewenvironment{numcases}[1]%
{$$\numc@opts \setbox\z@\hbox
 {\advance\c@equation\@ne\def\@currentlabel{\p@equation\theequation}% local
 $\displaystyle {#1\null}\m@th$}%
 \numc@setsub 
 \setbox\tw@\vbox\bgroup
 \stepcounter{equation}\def\@currentlabel{\p@equation\theequation}% 
 \global\@eqnswtrue\m@th \everycr{}\tabskip\numc@left\let\\\@eqncr
 \halign to\dimen@ii \bgroup \kern\wd\z@ \kern10
 %14
 \p@ % assume width of brace
 \tabskip\z@skip \global\@eqcnt\@ne $\displaystyle{##}$\hfil 
 &\global\@eqcnt\tw@ \skip@10\p@ \advance\skip@\tw@\arraycolsep \hskip\skip@
 ##\unskip\hfil\tabskip\@centering% \unskip removes space if no explanations
 &\global\@eqcnt\thr@@\hbox to\z@\bgroup\hss##\egroup\tabskip\z@skip\cr
}{\@@eqncr \egroup % end \halign, which does not contain first column or brace
 \global\advance\c@equation\m@ne 
%Measure the natural width of the alignment
 \unskip\unpenalty\unskip\unpenalty \setbox\z@\lastbox % grab last line
 \nointerlineskip \copy\z@ % then put it back
 \global\dimen@i\wd\z@ 
 \setbox\z@\hbox{\hskip-\numc@left\unhbox\z@}% Measure its natural width
 \ifdim \wd\z@<\dimen@i \global\dimen@i\wd\z@ \fi
\egroup% end \vbox (box\tw@, box\z@ is restored to LHS)
\hbox to\dimen@ii{\m@th % assemble the whole equation
 \hskip\numc@left 
 \hbox to\dimen@i{$\displaystyle \box\z@ % parameter #1
 \dimen@\ht\tw@ \advance\dimen@\dp\tw@ % get size of brace
 \left\{\vcenter to\dimen@{\vfil}\right.\n@space % make brace
 $\hfil}\hskip\@centering % finished first part (filled whole line)
 \kern-\dimen@ii % backspace the full width
 $\vcenter{\box\tw@}$% overlay the alignment
 }% end the \hbox to\dimen@ii
\numc@resetsub
$$\global\@ignoretrue}

\newcommand{\rightharpoonupfill}{$\m@th \mathord-\mkern-6mu \cleaders\hbox{$\mkern-2mu \mathord- \mkern-2mu$}\hfill\mkern-6mu \mathord \rightharpoonup$}
\makeatother

\newdimen\lengtharrow
\newbox\exponentbox

\newcommand{\Rightharpoonupts}[1]%
{\setbox\exponentbox=\hbox{$#1$}
\lengtharrow=\wd\exponentbox
\ifdim \lengtharrow=0pt
		\mathrel{\smash{\mathop{\hbox to 6mm{\rightharpoonupfill}}}}
\else \ifdim\lengtharrow<6mm
\lengtharrow=6mm
\else \advance\lengtharrow by 1mm
\fi
		\overset{\raisebox{3pt}{$\scriptstyle#1$}}{\mathrel{\smash{\hbox to\lengtharrow{\rightharpoonupfill}}}}
	\fi}
	
\newcommand{\Rightharpoonupds}[1]%
{\setbox\exponentbox=\hbox{$#1$}
\lengtharrow=\wd\exponentbox
\ifdim \lengtharrow=0pt
		\mathrel{\smash{\mathop{\hbox to 6mm{\rightharpoonupfill}}}}
\else \ifdim\lengtharrow<6mm
\lengtharrow=6mm
\else \advance\lengtharrow by 1mm
\fi
		\overset{\raisebox{5pt}{$\scriptstyle#1$}}{\mathrel{\smash{\hbox to\lengtharrow{\rightharpoonupfill}}}}
	\fi}
	
\def\Rightharpoonup#1{\mathchoice{\Rightharpoonupds{#1}}{\Rightharpoonupts{#1}}{\Rightharpoonupts{#1}}{\Rightharpoonupts{#1}}}

\newcommand{\longrightharpoonup}{\Rightharpoonup{}}



\multlinegap0pt
\newenvironment{enumeratei}
{\bgroup\def\theenumi{\roman{enumi}}\def\theenumii{\arabic{enumii}}\begin{enumerate}}
{\end{enumerate}\egroup}
\newcommand{\RedefinitSymbole}[1]{%
\expandafter\let\csname old\string#1\endcsname=#1
\let#1=\relax
\newcommand{#1}{\csname old\string#1\endcsname\,}%
}
\RedefinitSymbole{\forall} \RedefinitSymbole{\exists}

\newcommand{\<}{\bigl\langle}
\renewcommand{\>}{\bigr\rangle}
\newcommand{\Psfrac}[2]{(\sfrac{#1}{#2})}
\newcommand{\psfrac}[2]{\sfrac{(#1)}{#2}}
\newcommand{\spfrac}[2]{\sfrac{#1}{(#2)}}
\newcommand{\pspfrac}[2]{\sfrac{(#1)}{(#2)}}

\DeclareMathOperator{\meas}{meas}
\DeclareMathOperator{\spt}{spt}

\newcommand{\intTTT}{\int_0^T\hspace*{-2.5mm}\int_{\TT}}

\newcommand{\labeleq}[1]{}
\newenvironment{subnumcasess}{\begin{equation*}\begin{cases}}{\end{cases}\end{equation*}}

\newtheorem{thm}{Theorem}[section]
\newtheorem{prop}[thm]{Proposition}
\newtheorem{lem}[thm]{Lemma}

\theoremstyle{definition}
\newtheorem{definition}[thm]{Definition}
\newtheorem{assumption}[thm]{Assumption}

\theoremstyle{remark}
\newtheorem{remark}{Remark}

\newcommand{\TT}{\mathbb{T}^3}

\DeclareMathOperator{\curl}{curl}
\DeclareMathOperator{\Div}{div}
\newcommand{\R}{\mathcal{R}}
\newcommand{\di}{\,\mathrm{d}}
\DeclareMathOperator{\D}{D}
\newcommand{\Dt}{\dfrac{\mathrm{d}}{\mathrm{d} t}}
\let\ov\overline

\let\ep\varepsilon
\newcommand{\phil}{\phi_\lambda}
\newcommand{\phip}{\phi_p}

\usepackage{xfrac,xcolor,graphicx}\graphicspath{{./figures/}}\providecommand{\skpt}{\par\smallskip\noindent}\providecommand{\pointir}{.\ }\providecommand{\eg}{e.g.}\providecommand{\ie}{i.e.}\providecommand{\resp}{resp.}\providecommand{\cf}{cf.}\providecommand{\etc}{etc.}\let\Subsection\subsection\let\Subsubsection\subsubsection
